%% Security Analysis Extension for OrbitalChain
%% Add this content to Section 5.1 of your manuscript

\subsection{Extended Formal Security Analysis}

Building upon the component-level security established in Theorems 5.1 and 5.2, we now provide a comprehensive protocol-level security analysis using the simulation-based security paradigm.

\subsubsection{Ideal Functionality Definition}

We first formally define the ideal functionality that captures the desired security properties of our privacy-preserving streaming truth discovery protocol.

\begin{definition}[Ideal Functionality $\mathcal{F}_{\text{truth}}$]
\label{def:ideal_functionality}
The ideal functionality $\mathcal{F}_{\text{truth}}$ interacts with $K$ data providers $\{DP_1, \ldots, DP_K\}$ and $n$ satellites $\{S_1, \ldots, S_n\}$ as follows:

\textbf{Initialization:}
\begin{itemize}
    \item Set epoch counter $l \leftarrow 0$
    \item Initialize weights $w_k \leftarrow 1$ for all $k \in [K]$
    \item Initialize accumulated distances $st(k) \leftarrow 0$ for all $k \in [K]$
\end{itemize}

\textbf{Truth Discovery Phase (per epoch $l$):}
\begin{enumerate}
    \item \textbf{Input:} Receive sensing value $x_l^k$ from each $DP_k$
    \item \textbf{Computation:}
    \begin{itemize}
        \item Compute weighted sum: $t \leftarrow \sum_{k=1}^{K} w_k \cdot x_l^k$
        \item Compute weight sum: $z \leftarrow \sum_{k=1}^{K} w_k$
        \item Compute truth: $\Phi_l \leftarrow t / z$
    \end{itemize}
    \item \textbf{Output:} Send $\Phi_l$ to all satellites $\{S_1, \ldots, S_n\}$
\end{enumerate}

\textbf{Weight Update Phase (per epoch $l$):}
\begin{enumerate}
    \item \textbf{Computation:} For each $DP_k$:
    \begin{itemize}
        \item Compute squared error: $d_k \leftarrow (x_l^k - \Phi_l)^2$
        \item Update state: $st(k) \leftarrow st(k) \cdot \lambda + d_k$
        \item Compute total state: $st^* \leftarrow \sum_{k=1}^{K} st(k)$
        \item Update weight: $w_k \leftarrow -\log(st(k) / st^*)$
    \end{itemize}
    \item \textbf{Output:} Store updated weights for next epoch
\end{enumerate}
\end{definition}

\subsubsection{Security Model and Adversarial Assumptions}

\begin{definition}[Semi-Honest Adversary Model]
\label{def:adversary_model}
We consider a probabilistic polynomial-time (PPT) adversary $\mathcal{A}$ that:
\begin{itemize}
    \item Controls up to $t < n$ satellites (where $n$ is the total number of satellites in a shard)
    \item Follows the protocol specification honestly but attempts to learn additional information from the protocol transcript
    \item Cannot collude across different shards due to geographic separation
    \item Has access to all internal states of corrupted parties
\end{itemize}
\end{definition}

\begin{definition}[View of a Party]
\label{def:view}
The view of satellite $S_i$ during protocol execution $\pi$, denoted $\text{VIEW}_i^\pi$, consists of:
\begin{equation}
\text{VIEW}_i^\pi = \left( \{[x_l^k]_i\}_{k \in [K]}, \{[Tri_0]_i, [Tri_1]_i\}_{k \in [K]}, \{(u_j, v_j)\}_{j}, \{GC_{\text{output}}\}, r_i \right)
\end{equation}
where $r_i$ represents the random coins used by $S_i$.
\end{definition}

\subsubsection{Main Security Theorem}

\begin{theorem}[Protocol Security]
\label{thm:protocol_security}
The privacy-preserving streaming truth discovery protocol $\pi_{\text{truth}}$ (Algorithm 2) securely realizes the ideal functionality $\mathcal{F}_{\text{truth}}$ in the semi-honest model, under the following assumptions:
\begin{enumerate}
    \item The additive secret sharing scheme is perfectly secure
    \item Beaver's multiplication triple protocol is secure against semi-honest adversaries
    \item The garbled circuit construction satisfies simulation-based security
\end{enumerate}
\end{theorem}

\begin{proof}
We prove security by constructing a simulator $\mathcal{S}$ that, given only access to the ideal functionality $\mathcal{F}_{\text{truth}}$ and the set of corrupted parties, produces a view computationally indistinguishable from the real protocol execution.

\textbf{Simulator Construction:}

Let $\mathcal{C} \subset \{S_1, \ldots, S_n\}$ denote the set of corrupted satellites with $|\mathcal{C}| = t < n$. The simulator $\mathcal{S}$ operates as follows:

\textbf{Step 1: Simulating Secret Shares}

For each epoch $l$ and each data provider $DP_k$:
\begin{itemize}
    \item $\mathcal{S}$ receives the final output $\Phi_l$ from $\mathcal{F}_{\text{truth}}$
    \item For each corrupted satellite $S_i \in \mathcal{C}$, $\mathcal{S}$ generates simulated shares:
    \begin{equation}
    [\tilde{x}_l^k]_i \xleftarrow{\$} \mathbb{F}_q
    \end{equation}
    \item These shares are uniformly random and independent of the actual sensing values
\end{itemize}

\textbf{Step 2: Simulating Multiplication Triples}

For each required multiplication operation:
\begin{itemize}
    \item $\mathcal{S}$ generates random triple shares for corrupted parties:
    \begin{equation}
    ([\tilde{a}]_i, [\tilde{b}]_i, [\tilde{c}]_i) \xleftarrow{\$} \mathbb{F}_q^3 \text{ for } S_i \in \mathcal{C}
    \end{equation}
    \item When simulating Beaver's protocol, $\mathcal{S}$ computes:
    \begin{align}
    \tilde{u} &\xleftarrow{\$} \mathbb{F}_q \\
    \tilde{v} &\xleftarrow{\$} \mathbb{F}_q
    \end{align}
    \item $\mathcal{S}$ broadcasts $(\tilde{u}, \tilde{v})$ as the masked differences
\end{itemize}

\textbf{Step 3: Simulating Garbled Circuits}

For secure division ($GC_{\text{div}}$) and logarithm ($GC_{\text{div+log}}$) operations:
\begin{itemize}
    \item Let $m$ be the expected output (known from $\mathcal{F}_{\text{truth}}$)
    \item $\mathcal{S}$ invokes the garbled circuit simulator $\mathcal{S}_{GC}$ (which exists by the security of Yao's protocol \cite{lindell2009proof}):
    \begin{equation}
    \widetilde{GC}, \{\tilde{l}_i\} \leftarrow \mathcal{S}_{GC}(1^\kappa, m)
    \end{equation}
    \item $\mathcal{S}$ provides the simulated garbled circuit and labels to corrupted parties
\end{itemize}

\textbf{Indistinguishability Argument:}

We show that $\text{VIEW}_{\mathcal{C}}^{\pi} \approx_c \text{VIEW}_{\mathcal{C}}^{\mathcal{S}}$ through a hybrid argument.

\textbf{Hybrid $H_0$:} Real protocol execution.

\textbf{Hybrid $H_1$:} Replace real secret shares with random shares.
\begin{itemize}
    \item By the perfect secrecy of additive secret sharing (Theorem 5.1), for any $t < n$ shares, the conditional distribution satisfies:
    \begin{equation}
    H(x_l^k | \{[x_l^k]_i\}_{i \in \mathcal{C}}) = H(x_l^k)
    \end{equation}
    \item Therefore, $H_0 \equiv H_1$ (perfectly indistinguishable)
\end{itemize}

\textbf{Hybrid $H_2$:} Replace real Beaver protocol messages with simulated ones.
\begin{itemize}
    \item In the real protocol, the broadcast values are:
    \begin{align}
    u &= x - a \\
    v &= y - b
    \end{align}
    where $(a, b)$ are uniformly random in $\mathbb{F}_q$
    \item Since $a, b$ are independent of $x, y$, the values $u, v$ are uniformly distributed
    \item By Theorem 5.2, $H_1 \equiv H_2$ (perfectly indistinguishable)
\end{itemize}

\textbf{Hybrid $H_3$:} Replace real garbled circuits with simulated ones.
\begin{itemize}
    \item By the simulation security of garbled circuits \cite{lindell2009proof}, there exists a PPT simulator $\mathcal{S}_{GC}$ such that:
    \begin{equation}
    \{GC, \{l_i\}_{i \in \text{input}}\} \approx_c \mathcal{S}_{GC}(1^\kappa, f(x))
    \end{equation}
    \item Therefore, $H_2 \approx_c H_3$
\end{itemize}

Since $H_3$ is exactly the simulated view, we conclude:
\begin{equation}
\text{VIEW}_{\mathcal{C}}^{\pi} \equiv H_0 \equiv H_1 \equiv H_2 \approx_c H_3 = \text{VIEW}_{\mathcal{C}}^{\mathcal{S}}
\end{equation}

This completes the proof. \qed
\end{proof}

\subsubsection{Composition Security}

\begin{theorem}[Secure Composition with D-Stream]
\label{thm:composition}
The composition of the privacy-preserving streaming truth discovery protocol (Algorithm 2) with the D-Stream clustering algorithm (Algorithm 1) preserves the privacy of individual sensing values. Formally, the composed protocol $\pi_{\text{composed}} = \pi_{\text{truth}} \circ \pi_{\text{D-Stream}}$ reveals no information about individual inputs $\{x_l^k\}$ beyond what is revealed by the legitimate outputs $\{\Phi_l, w_k\}$.
\end{theorem}

\begin{proof}
We prove this theorem by analyzing the information flow between the two protocols.

\textbf{Information Flow Analysis:}

The D-Stream algorithm (Algorithm 1) receives the following inputs from the truth discovery protocol:
\begin{enumerate}
    \item Discovered truth values $\{\Phi_l\}_{l \geq 1}$
    \item Updated satellite weights $\{w_k\}_{k \in [K]}$
    \item Data point coordinates derived from aggregated truths
\end{enumerate}

\textbf{Claim 1:} The D-Stream algorithm never accesses secret-shared values.

\textit{Proof of Claim 1:} By inspection of Algorithm 1, all operations are performed on:
\begin{itemize}
    \item Data points $x = (x_1, \ldots, x_d)$ which represent aggregated/public information
    \item Weights $w$ which are the output of Algorithm 2
    \item Density grid statistics computed from public truths
\end{itemize}
No line in Algorithm 1 references $[x_l^k]_i$ or any intermediate MPC values. \qed

\textbf{Claim 2:} The outputs of Algorithm 2 constitute a valid leakage function.

\textit{Proof of Claim 2:} The ideal functionality $\mathcal{F}_{\text{truth}}$ explicitly defines $\{\Phi_l, w_k\}$ as legitimate outputs. Any function of these outputs is therefore permitted leakage. The D-Stream clustering produces clusters based solely on these permitted outputs. \qed

\textbf{Formal Composition Argument:}

Let $\mathcal{L}(\{x_l^k\}) = \{\Phi_l, w_k\}$ denote the leakage function. By the Universal Composability (UC) framework \cite{canetti2000security}:

\begin{equation}
\pi_{\text{composed}} \text{ securely realizes } \mathcal{F}_{\text{composed}} \Leftrightarrow 
\begin{cases}
\pi_{\text{truth}} \text{ securely realizes } \mathcal{F}_{\text{truth}} \\
\pi_{\text{D-Stream}}(\mathcal{L}(\cdot)) \text{ uses only } \mathcal{L}(\cdot)
\end{cases}
\end{equation}

Since both conditions are satisfied (by Theorem \ref{thm:protocol_security} and Claim 1), the composition is secure. \qed
\end{proof}

\subsubsection{Security Against Active Adversaries}

While our main analysis focuses on the semi-honest model, we briefly discuss extensions to handle active adversaries within the PBFT consensus framework.

\begin{theorem}[Byzantine Fault Tolerance]
\label{thm:bft}
The OrbitalChain consensus mechanism tolerates up to $f < n/3$ Byzantine satellites per shard while maintaining:
\begin{enumerate}
    \item \textbf{Safety:} All honest satellites agree on the same truth values
    \item \textbf{Liveness:} The protocol eventually produces outputs if fewer than $f$ satellites are faulty
\end{enumerate}
\end{theorem}

\begin{proof}
The proof follows from the standard PBFT security analysis \cite{castro1999practical} adapted to our shard structure.

\textbf{Safety:} Suppose two honest satellites $S_i$ and $S_j$ commit to different truth values $\Phi$ and $\Phi'$ at epoch $l$. This requires:
\begin{itemize}
    \item $S_i$ received $2f+1$ COMMIT messages for $\Phi$
    \item $S_j$ received $2f+1$ COMMIT messages for $\Phi'$
\end{itemize}

Since $n = 3f+1$, the intersection of these sets contains at least:
\begin{equation}
(2f+1) + (2f+1) - n = f+1
\end{equation}
satellites. At least one of these must be honest (since at most $f$ are Byzantine), contradicting the assumption that an honest satellite sent COMMIT for both values.

\textbf{Liveness:} With at most $f$ Byzantine satellites, at least $2f+1$ honest satellites remain. The leader rotation mechanism ensures that an honest leader is eventually selected, who will drive the protocol to completion. \qed
\end{proof}

\subsubsection{Complexity Analysis}

\begin{theorem}[Communication Complexity]
\label{thm:complexity}
The communication complexity of Algorithm 2 per epoch is:
\begin{equation}
O(K \cdot n \cdot \kappa + n^2 \cdot |GC|)
\end{equation}
where $K$ is the number of data providers, $n$ is the number of satellites, $\kappa$ is the security parameter, and $|GC|$ is the garbled circuit size.
\end{theorem}

\begin{proof}
\textbf{Secret Sharing Phase:}
\begin{itemize}
    \item Each DP sends $n$ shares of size $O(\kappa)$ bits
    \item Total: $O(K \cdot n \cdot \kappa)$
\end{itemize}

\textbf{Beaver's Multiplication:}
\begin{itemize}
    \item Each multiplication requires broadcasting $(u, v)$: $O(n)$ messages of $O(\kappa)$ bits
    \item Number of multiplications per epoch: $O(K)$
    \item Total: $O(K \cdot n \cdot \kappa)$
\end{itemize}

\textbf{Garbled Circuits:}
\begin{itemize}
    \item Division circuit: $O(|GC_{\text{div}}|)$ 
    \item Division + Logarithm circuit: $O(|GC_{\text{div+log}}|)$
    \item Oblivious transfer: $O(n \cdot \kappa)$
    \item Total: $O(n^2 \cdot |GC|)$ considering all pairwise interactions
\end{itemize}

Combining all phases yields the stated complexity. \qed
\end{proof}


%% References to add
\begin{thebibliography}{99}

\bibitem{canetti2000security}
R. Canetti, ``Security and composition of multiparty cryptographic protocols,'' \textit{Journal of Cryptology}, vol. 13, pp. 143--202, 2000.

\bibitem{lindell2009proof}
Y. Lindell and B. Pinkas, ``A proof of security of Yao's protocol for two-party computation,'' \textit{Journal of Cryptology}, vol. 22, pp. 161--188, 2009.

\bibitem{castro1999practical}
M. Castro and B. Liskov, ``Practical Byzantine fault tolerance,'' \textit{Third Symposium on Operating Systems Design and Implementation}, pp. 359--368, 1999.

\bibitem{goldreich2009foundations}
O. Goldreich, \textit{Foundations of Cryptography: Volume 2, Basic Applications}. Cambridge University Press, 2009.

\end{thebibliography}
